\begin{scanpage}{0052}{35}
and are conical surfaces, the generatrices of which form the angle $\pi/2+\gamma$ with the positive direction of the $z$ axis. Thus, with oblique irradiation of the wedge by a plane wave, the field created by the nonuniform part of the current is a set of conical waves diverging from the wedge edge. The normals to the phase surfaces of these waves form an angle $\gamma$ with the positive direction of the $z$ axis and are shown in Figure 15. These waves may be represented in a more graphic form if one introduces the components (see Figure 15):
\[\left.\begin{aligned}E_{\gamma}&=E_r\cos\gamma-E_z\sin\gamma,\\H_{\gamma}&=H_r\cos\gamma-H_z\sin\gamma.\end{aligned}\right\}\tag{5.15}\]
Then the final expressions for the fringing field in the far zone will have the form
\[\left.\begin{aligned}E_{\gamma}&=H_\phi=-\frac1{\sin\gamma}E_z,\\H_{\gamma}&=-E_\phi=-\frac1{\sin\gamma}H_z.\end{aligned}\right\}\tag{5.16}\]
Now we are able to proceed to the application of the results which have been obtained for the solution of specific diffraction problems. The simplest of them is the problem of diffraction by an infinitely long strip which has a rigorous solution [23] in the form of Mathieu function series. However, in the quasi-optical region when the width of the strip is large in comparison with the wavelength, these series have a poor convergence and are not suitable for numerical calculations. Therefore, the requirement arises for approximation equations which are useful in the quasi-optical region. The derivation of such equations for a field scattered by a strip will be given in the following section.
\sectiontitle{6}{Diffraction by a Strip}
Let us investigate diffraction by an infinitely thin, ideally conducting strip which has a width of $2a$ and an unlimited length. The orientation of the strip in space is shown in Figure 16.
\end{scanpage}
